% ============================================================
% Section 2: Related Work
% Negative diagnostic study: closest-work positioning
% ============================================================

\section{Related Work}
\label{sec:related}

The study sits at the intersection of graph-characterization metrics, heterophilous node classification, graph learning model selection, and tests of whether graph structure is necessary. Its contribution is a frozen decision evaluation of simple diagnostics, not a new homophily measure, graph architecture, or general model-selection benchmark.

\subsection{Graph Characteristics and Relative Model Performance}

Early homophily measures summarize whether adjacent nodes share labels, but a single global fraction omits class imbalance, feature geometry, neighborhood distributions, and local variation~\citep{mcpherson2001homophily,zhu2020beyond}. Subsequent work introduced post-aggregation similarity, adjusted homophily, label informativeness, and local or class-conditional characterizations to better describe when message passing helps~\citep{luan2022linkx,platonov2023characterizing,loveland2024local}. Accordingly, our contribution is not the established observation that homophily is insufficient~\citep{ma2021homophily}.

The closest diagnostic work goes further. \citet{luan2023whendo} analyze node distinguishability and propose feature-based hypothesis-testing metrics with thresholds intended to reveal the advantage of graph-aware models. \citet{zheng2024trihom} disentangle label, structural, and feature homophily; their 31-dataset study compares Tri-Hom with 17 metrics, including correlations with GCN-, GraphSAGE-, and GAT-versus-MLP performance gaps. Under a heterophilous stochastic block model, \citet{pmlr-v235-wang24u} relate separability gain to neighborhood-distribution distance and degree rather than to global homophily alone. These studies motivate richer diagnostics but primarily evaluate theoretical quantities or performance association. We instead freeze specified statistics as graph/MLP/abstain policies and evaluate the cost of their actions through regret and coverage.

Benchmark design also narrows the novelty claim. Repeated-split evaluation and standardized graph collections expose conclusions that depend on data splits, tuning, or dataset choice~\citep{shchur2018pitfalls,hu2020ogb}. \citet{platonov2023heterophilous} reassess whether reported progress on heterophilous benchmarks survives stronger protocols, while \citet{luan2024reevaluating} quantitatively re-evaluate both models and homophily metrics. Comparisons across datasets further require treating datasets, rather than pooled seeds, as the inferential units~\citep{demsar2006statistical}. We therefore do not claim the first negative evaluation of homophily-related diagnostics. Our narrower question is whether transparent rules whose label-dependent graph statistics use training labels only retain incremental decision value over matched trivial policies on real dataset--seed units when negative outcomes are preserved without threshold redesign.

\subsection{Heterophily-Aware Architectures}

Heterophily-aware models change which signals are propagated or how they are combined. H2GCN separates ego and neighbor representations and uses two-hop neighborhoods~\citep{zhu2020beyond}; LINKX separates structural and feature processing~\citep{lim2021linkx}; MixHop and GPR-GNN combine multiple propagation depths~\citep{abu2019mixhop,chien2021gprgnn}; FAGCN, GloGNN, and ACM-GCN adapt filters or communication patterns~\citep{fagcn2021,li2022glognn,luan2022linkx}. Geom-GCN constructs geometric neighborhoods, while high-frequency and flexible spectral designs broaden the signals a graph model can retain~\citep{pei2020geomgcn,he2022bernnet,wang2022jacobi}. The wider node-classification design space also includes simplified propagation, deeper residual propagation, revised attention, and graph-augmented MLPs~\citep{wu2019sgc,chen2020gcnii,brody2022gatv2,zhang2021gamlp}. These architectures demonstrate that low one-hop homophily need not make graph structure useless. They also make architecture-specific prediction difficult: a statistic that describes one-hop mixing need not identify which portfolio wins after tuning. Our benchmark therefore places GCN, GAT, GraphSAGE, H2GCN, LINKX, and GPR-GNN in the graph portfolio and evaluates a portfolio action rather than presenting a new heterophily architecture.

Depth introduces a separate limitation: repeated propagation can erase discriminative variation or contract node representations~\citep{oono2020exponential,chen2020measuring}. DropEdge and PairNorm address this behavior through edge sampling and representation normalization~\citep{rong2020dropedge,zhao2020pairnorm}. These results reinforce the distinction between describing a graph and forecasting the outcome of a tuned portfolio: realized graph utility depends jointly on data, architecture, optimization, and regularization.

\subsection{Graph Learning Model Selection}

Graph model selection has been studied directly. Auto-GNN searches neural architectures for a given graph task~\citep{zhou2020autognn}. MetaGL learns to select graph-learning models on a new graph from meta-graph features and historical performance~\citep{park2023metagl}. GLEMOS scales this problem to 366 models and 457 graphs and benchmarks multiple instantaneous selection algorithms~\citep{park2023glemos}. Consequently, the present 11-dataset experiment is neither the first nor the largest graph-model-selection benchmark.

Our estimand differs from ranking many individual models across a large meta-dataset. We ask whether predeclared, interpretable statistics can support a binary graph-portfolio-versus-MLP-portfolio action without using test labels. Each architecture receives the same four-trial validation grid, but the total search budgets are not equal: the graph action selects among six architectures and the feature-only action contains one MLP. The final test outcome is evaluated once, and the diagnostic rules are compared with always-graph, always-MLP, random, and direct validation-selection policies on the same 110 units. This design treats diagnostics as deployable decisions and makes oracle-portfolio regret and abstention coverage explicit.

\subsection{When Is the Graph Necessary?}

Strong feature-only baselines are essential because some graph benchmarks can be solved largely from node features. Prior evaluations show that model rankings depend on data splits and training procedures~\citep{shchur2018pitfalls}. Well-tuned feature-only MLPs can also be competitive when node features already encode graph information~\citep{katsman2026necessary}. Other controls compare predictions with and without meaningful edges, showing that models can exploit graph structure when it hurts generalization~\citep{bechler2024graphs}. These findings motivate the graph-versus-MLP action but do not supply a cheap rule for making it.

Our intervention complements these controls. It holds every node's degree fixed while randomizing edge organization. It does not isolate homophily: class mixing, paths, components, clustering, and other higher-order properties can change together. We use the intervention only to reject the interpretation that diagnostic failure demonstrates graph irrelevance. The main contribution remains the decision audit: useful graph signal in the intervention does not establish that the evaluated diagnostic policies improve graph-versus-MLP selection.

\subsection{Scoped Mechanistic Context}

The fixed-degree Gaussian calculation in Appendix~\ref{sec:snr} is included to explain why correlation, degree, and feature strength can interact under a prescribed linear average. It is adjacent to contextual stochastic-block-model recovery theory and finite-depth embedding analyses~\citep{lu2023contextual,pmlr-v269-dalle25a}, but it is not a new recovery threshold, Bayes-error result, or trained-GNN guarantee. The negative benchmark does not depend on that calculation being a universal theory; it evaluates the diagnostics directly under the frozen empirical protocol.
